\documentclass[12pt,a4paper]{article}
\usepackage[margin=1in]{geometry}
\usepackage{amsmath,amssymb,amsthm,mathtools,booktabs,array,fontspec}
\usepackage[hidelinks]{hyperref}
\defaultfontfeatures{Ligatures=TeX}
\IfFontExistsTF{Times New Roman}{\setmainfont{Times New Roman}}{\setmainfont{TeX Gyre Termes}}
\newtheorem{theorem}{Theorem}[section]
\newtheorem{proposition}[theorem]{Proposition}
\newtheorem{lemma}[theorem]{Lemma}
\theoremstyle{definition}
\newtheorem{conjecture}[theorem]{Conjecture}
\newcommand{\tr}{\operatorname{tr}}
\newcommand{\ran}{\operatorname{ran}}
\newcommand{\E}{\mathbb E}
\hypersetup{pdftitle={No View from Nowhere: Read-out Closure and the Vacuum Curvature Test in Q-RAIF A},pdfauthor={Sidong Liu}}
\title{\textbf{No View from Nowhere}\\\large Read-out Closure and the Vacuum Curvature Test in Q-RAIF A}
\author{Sidong Liu, PhD\\\small iBioStratix Ltd\\\small\texttt{sidongliu@hotmail.com}}
\date{October 2026 working revision\\\small Series record DOI:
\href{https://doi.org/10.5281/zenodo.20284454}{10.5281/zenodo.20284454}}
\begin{document}
\emergencystretch=2em
\raggedbottom
\maketitle

\begin{abstract}
A source can be physical and still be missed by a restricted interface.
For a finite-dimensional readout channel $\Lambda$ and source observable
$H$, the operator-norm distance from $H$ to the dual range equals half
the largest energy difference between states with the same output. It is
also the smallest worst-case error of any energy estimate based on the
exact output state, including nonlinear estimates. An actual finite
instrument adds a separate sampling cost. We exhibit a positive two-part
source whose hidden interaction gives a nonzero error floor, then remove
that floor by adding a priced mixed-memory record of the same source.

Complete null stress data have a metric-trace blind direction. A real
vacuum-density difference therefore cannot be dismissed as a zero-point
convention merely because null readout misses it. A matched Gaussian
material supplies a complementary test: changing its contact changes the
fluctuation operator and its logarithmic vacuum coefficient, with one
physical calibration required across configurations.

These results turn readout visibility into a source-certification problem.
They support, without proving, a broader closure conjecture: a spacetime
description must specify or generate the interfaces that select its
records and sources. Whether those interfaces can be fully internalized
is an independent question. Neither interface blindness nor the closure
conjecture computes the cosmological constant or switches off physical
vacuum gravitation.
\end{abstract}

\section{What the reader leaves outside}
\label{sec:blade}

A laboratory description places its detectors inside spacetime and lets
their records describe the rest of the world. It also needs a rule for
what counts as a detector, a record and a relevant source. The question
behind ``no view from nowhere'' is whether those rules can themselves be
generated and represented inside the same description, with their physical
execution accounted for.

This is a question about structural interfaces, not a claim that a human
mind creates the world. Q-RAIF A obtains a Lorentzian propagation form
under a common quadratic-shell condition and its independence assumptions;
the dynamics need not select one readout cone \cite{Liu2026QRAIFA}. The
construction makes the selection problem concrete. It does not prove
that an observer-independent formulation is impossible. HAFF G's finite
programming obstruction likewise distinguishes a programmed target
operation from its joint executor; it does not prohibit stored
self-descriptions or typed fixed points.

Vacuum energy supplies a particularly demanding test. A description may
know how to read radiation while lacking a trace-sensitive source datum.
Calling the missing datum invisible does not settle its gravitational
effect. We first determine exactly what a finite interface can certify,
including how much source can remain hidden. We then apply that distinction
to null stress, physical vacuum differences and changes of apparatus. The
closure question returns after these concrete results.

\section{A source certificate for a finite interface}
\label{sec:source}

Let $\Lambda$ be a trace-preserving quantum channel between
finite-dimensional matrix algebras and $H=H^*$ a specified source
observable. The source datum is $\tr H\rho$. It may be an energy or a
bounded smeared source in a regulated model. Its definition precedes the
response calculation. The results below do not apply an operator-norm
bound to arbitrary unbounded continuum stress.

\begin{proposition}[Exact source completeness]
\label{prop:source-complete}
For all input states,
\begin{equation}
 \Lambda\rho=\Lambda\sigma\ \Longrightarrow\
 \tr H\rho=\tr H\sigma
\end{equation}
if and only if $H=\Lambda^*O$ for a Hermitian output observable $O$.
For such an $O$,
\begin{equation}
 |\tr H(\rho-\sigma)|
 \leq\inf_{a\in\mathbb R}\|O-aI\|_\infty
                  \|\Lambda\rho-\Lambda\sigma\|_1.
 \label{eq:source-stability}
\end{equation}
\end{proposition}
\begin{proof}
A Hermitian $D\in\ker\Lambda$ has trace zero. For sufficiently small
$t>0$, $I/d\pm tD$ are states. The implication thus gives $\tr HD=0$.
Finite-dimensional duality identifies the orthogonal complement of the
kernel with $\ran\Lambda^*$. Conversely, $H=\Lambda^*O$ makes its
expectation constant on every output fibre. The final inequality is
trace/operator-norm duality applied to the trace-zero output difference
\cite{Watrous2018}.
\end{proof}

Exact source completeness does not require full state tomography. When
it fails, the missing source has a sharp quantitative size:
\begin{equation}
 \delta_H(\Lambda)
 =\inf_{O=O^*}\|H-\Lambda^*O\|_\infty.
 \label{eq:delta}
\end{equation}

\begin{theorem}[Hidden-source size and minimax error]
\label{thm:minimax}
The infimum in \eqref{eq:delta} is attained, and
\begin{equation}
 \delta_H(\Lambda)
 =\frac12\max_{\Lambda\rho=\Lambda\sigma}
       |\tr H(\rho-\sigma)|
 =\inf_f\sup_\rho|\tr H\rho-f(\Lambda\rho)|.
 \label{eq:minimax}
\end{equation}
Here $f$ may be any real-valued, nonlinear function of the exact output
state. Thus no such function improves the worst-case absolute error of
an optimal affine estimate $f(\tau)=\tr O\tau$.
\end{theorem}

\begin{proof}
Work in the real vector space of Hermitian matrices. The range of
$\Lambda^*$ is a finite-dimensional closed subspace, so a best
operator-norm approximation to $H$ exists. Norm duality and the
Hahn--Banach distance formula give
\begin{equation}
 \delta_H(\Lambda)
 =\max\{\tr HD:D=D^*,\ \Lambda D=0,\ \|D\|_1\leq1\}.
 \label{eq:dual}
\end{equation}
The kernel is invariant under sign reversal, so the maximum could also
use an absolute value. Write $D=D_+-D_-$ and
$c=\tr D_+=\tr D_-\leq1/2$. For any state $\tau$, the matrices
\begin{equation}
 \rho=2D_++(1-2c)\tau,\qquad
 \sigma=2D_-+(1-2c)\tau
\end{equation}
are states with the same output and $\rho-\sigma=2D$. Conversely every
equal-output state pair gives an admissible $D=(\rho-\sigma)/2$.
This proves the first equality in \eqref{eq:minimax}, with attainment.

On a common output an estimate $f$ gives the same number for both states.
At least one of their absolute errors is at least half their source
difference. This proves the minimax lower bound. A minimizing $O$ in
\eqref{eq:delta} gives error at most $\delta_H$ on every state, proving
the upper bound. Trace preservation gives $\Lambda^*I=I$, so the
approximation space includes constant offsets as well as the nonconstant
observable components. No missing centering assumption is used.
\end{proof}

The proof uses standard quotient-norm duality; its consequence here is a
sharp source uncertainty for a specified physical interface. The theorem
concerns information in an \emph{exact output state}. It
does not give a finite observer that state for free. It also does not
grant a measurement of an arbitrary element of $\ran\Lambda^*$.
An actual instrument is needed for the next step.

\begin{proposition}[Finite-instrument source interval]
\label{prop:finite}
Suppose an allowed output POVM has effects $F_y$ and real scores $c_y$,
with $O=\sum_y c_yF_y$, score width
$w=\max_y c_y-\min_y c_y$, and
$\delta_O=\|H-\Lambda^*O\|_\infty$. Use this same instrument on $n$
independently acquired outputs of the same input state. With probability
at least $1-\zeta$, the sample mean $\overline c$ satisfies
\begin{equation}
 |\tr H\rho-\overline c|
 \leq\delta_O+w\sqrt{\frac{\log(2/\zeta)}{2n}}.
 \label{eq:certificate}
\end{equation}
\end{proposition}
\begin{proof}
The output score has mean $\tr O\Lambda\rho$. Its bias relative to the
source is at most $\delta_O$. Hoeffding's bounded-variable inequality
\cite{Hoeffding1963} bounds the deviation of its sample mean by the
second term: the centered score moment-generating function is at most
$\exp(\lambda^2w^2/8)$, and the independent product followed by the
Chernoff bound gives $2\exp(-2nt^2/w^2)$ for a two-sided deviation $t$.
Adding the bias proves \eqref{eq:certificate}; for $w=0$ the sampling
term vanishes directly.
\end{proof}

The interval is conditional on the source model, the actual effects and
their stable calibration. Copy preparation, readout, storage and the
independence of the acquired outputs are inputs, not consequences of
dual-range membership. Independently producing the copies is not cloning
an arbitrary unknown state. Existing references and spent ports remain
part of the physical ledger, even if their measurement is outside the
declared permission. The certificate may be used in a supplied response
law: if a response is $L$-Lipschitz in this source datum, its corresponding
error is bounded by $L$ times \eqref{eq:certificate}. That step does not
derive the response law itself.

\section{A hidden interaction, then its actual completion}
\label{sec:completion}

The distinction can be realized on one finite source rather than by
comparing unrelated channels. Take two source qubits $A,B$, with
\begin{equation}
 H=E_0I+\omega P_{1,A}+g Z_AZ_B,
 \qquad \omega\geq0,\quad E_0\geq|g|.
 \label{eq:ABsource}
\end{equation}
This is a positive energy observable, including the interaction. If only
$A$ is available, $\Lambda=\tr_B$ and the output observable
$O_A=E_0I+\omega P_{1,A}$ gives residual norm $|g|$. The two states
$|00\rangle$ and $|01\rangle$ have the same $A$ output and energy
difference $2g$. Theorem~\ref{thm:minimax} therefore gives
\begin{equation}
 \delta_H(\tr_B)=|g|.
 \label{eq:ABdelta}
\end{equation}
Even complete knowledge of the $A$ state leaves this error floor. It is
not a lack of data analysis, and it does not remove $gZ_AZ_B$ from $H$.

For an actual record use the partially mixed ports of HAFF E and F
\cite{Liu2026PaperE,Liu2026PaperF}. A fresh port initially in
$\eta_v=(I+vZ)/2$, $0<v\leq1$, is coupled by controlled $X$ from its
source qubit and read in the port's $Z$ basis. With source eigenvalue
$z=\pm1$, its output obeys
\begin{equation}
 \Pr(Y=y\mid z)=\frac{1+yvz}{2},\qquad
 F_y=\frac{I+yvZ}{2}.
 \label{eq:port}
\end{equation}
The coupling and subsequent readout give this law on the actual joint
state. They dephase source coherence in the measured basis; they are not
the L\"uders instrument of the unsharp effect. The source populations,
and therefore the expectation of the diagonal Hamiltonian
\eqref{eq:ABsource}, are unchanged.

With only the $A$ port, score the record by
\begin{equation}
 c_A(y)=E_0+\frac\omega2(1-y/v).
\end{equation}
Its mean is $\tr O_A\rho_A$, its width is $\omega/v$, and the residual
bias is bounded by $|g|$. Thus the same hidden-source floor is accompanied
by an attainable finite certificate, not just an abstract optimal
observable.

Now add a second independent port coupled to $B$ on each \emph{same}
source pair. The physical joint effects are $F_{y_A}\otimes F_{y_B}$;
the outcomes are independent conditional on the two source labels, not
independent after averaging over an arbitrary correlated source. Score
the paired record by
\begin{equation}
 c(y_A,y_B)=E_0+\frac\omega2(1-y_A/v)+g\frac{y_Ay_B}{v^2}.
 \label{eq:fullscore}
\end{equation}
Directly summing the four effects with these scores gives $H$. The
interaction bias is now zero. A score-width bound is
$w\leq\omega/v+2|g|/v^2$, so Proposition~\ref{prop:finite} gives a
finite full-source certificate on every two-qubit state. The extra port
has supplied the missing relation $Z_AZ_B$, not merely refined the old
$A$ marginal.

For $E_0=1$, $\omega=2$, $g=1/4$, $v=0.6$, a $95\%$ certificate of
absolute error at most $0.5$ is guaranteed by $328$ independent source
pairs and $328$ ports in the one-port protocol, or by $165$ source pairs
and $330$ ports in the completed protocol. These conservative budgets
use \eqref{eq:certificate}. They show a tradeoff between source acquisition
and port stock, not a universal saving in total energy. Each port brings
purity deficit $\log2-h((1+v)/2)$, where $h$ is binary entropy. The
known $v$, controlled gates, independent production, readout, paired
labels and retained archive all count. No thermal reset or perpetual
copy supply has been implemented by this example.

The source-completion problem is therefore constructive: identify a
hidden source direction, add a physical read of the missing relation,
and recalculate both the residual and the statistical budget. The
algebraic criterion alone would miss the costs and gains of that change.

\section{Algebraic completeness can still be unusable}
\label{sec:conditioning}

The opposite limit matters as well. Let a qubit dephasing channel obey
$\Lambda_\eta(I,X,Y,Z)=(I,\eta X,\eta Y,Z)$, $0\leq\eta\leq1$, and
take the positive observable $H=E(I+X)$, $E>0$. For $\eta>0$ the dual
range is full, $\delta_H=0$ and $O=E(I+X/\eta)$ reconstructs the energy.
The coefficient in \eqref{eq:source-stability} is $E/\eta$.
At $\eta=0$, the states $\rho_\pm=(I\pm X)/2$ have identical outputs
and energy difference $2E$, so $\delta_H=E$.

For $0<\eta<1$ these two outputs commute. Their relative entropy is
\begin{equation}
 d_\eta=\eta\log\frac{1+\eta}{1-\eta}
        =2\eta^2+O(\eta^4).
\end{equation}
With $n$ independent copies it is $nd_\eta$. The standard
Bretagnolle--Huber testing bound \cite{BretagnolleHuber1978} gives
$P_e\geq\tfrac14e^{-nd_\eta}$ for equal prior probabilities. For the
commuting distributions, Jensen bounds their squared Hellinger affinity
below by $e^{-nd_\eta}$, and Cauchy--Schwarz bounds total variation
above by $\sqrt{1-e^{-nd_\eta}}$, giving this testing bound. Thus
$P_e\leq\varepsilon<1/4$ requires
$n\geq\log(1/(4\varepsilon))/d_\eta$. A joint quantum measurement
cannot evade the classical bound for these commuting states. Measuring
$X$ attains an actual informative binary experiment if that measurement
is permitted; an additional noisy port introduces its own gain and cost.
Small output gain can make source resolution expensive long before the
exact range loses a direction. This is a finite-source example, not a
model of an absolute vacuum-density shift.

\section{The vacuum curvature test}
\label{sec:vacuum}

\subsection{Null data miss a physical trace}

\begin{lemma}[Complete null-stress kernel]
\label{lem:null-trace}
In a Lorentzian vector space of dimension $D\geq2$, a symmetric tensor
$S$ has $S_{\mu\nu}k^\mu k^\nu=0$ for every null $k$ if and only if
$S=\alpha g$.
\end{lemma}
\begin{proof}
In an orthonormal frame set $k=(1,\mathbf n)$, $|\mathbf n|=1$.
Comparing $\mathbf n$ and $-\mathbf n$ gives $S_{0i}=0$. The spatial
quadratic form is constant on the unit sphere, so
$S_{ij}=-S_{00}\delta_{ij}$. Hence $S=-S_{00}g$. The two null
directions give the same result for $D=2$; the converse is immediate.
\end{proof}

A physical vacuum-density difference
$-\delta\rho\,g_{\mu\nu}$ lies in this kernel. Null tomography, even
ideal and complete in its own menu, cannot certify that source component.
If two smooth stress tensors are separately conserved on a connected
domain, their blind difference $\alpha(x)g_{\mu\nu}$ must satisfy
$\nabla^\mu(\alpha g_{\mu\nu})=\partial_\nu\alpha=0$.
The conserved ambiguity is therefore a constant metric-trace term,
not an arbitrary position-dependent vacuum profile.
The closed task in the null-cluster paper \cite{Liu2026NullCluster}
illustrates both sides: null records estimate a gradient contribution
efficiently, yet remain exactly blind to a potential-energy difference.
In HAFF D's supplied JT response, the null components of a constant
vacuum shift vanish while the dilaton shifts by
$\kappa L^2\delta\rho$ \cite{Liu2026HAFFD}. Such a shift gives a
response uncertainty if its source is unmeasured; it is not evidence for
degravitation. A trace-sensitive datum or independently fixed parameter
is required.

\subsection{A convention is not a changed stress tensor}

The semiclassical source equation is
\begin{equation}
 G_{\mu\nu}+\Lambda g_{\mu\nu}
       =8\pi G\langle T_{\mu\nu}\rangle_{\rm ren}.
\end{equation}
A Lorentz-invariant vacuum contribution has the form
$\langle T_{\mu\nu}\rangle_{\rm vac}=-\rho_{\rm vac}g_{\mu\nu}$.
The renormalized parameter and its common calibration matter; a bare
regulator sum $\tfrac12\sum_{\mathbf k}\hbar\omega_{\mathbf k}$ is not
a parameter-free prediction of that stress \cite{Weinberg1989,Wald1977}.

Two changes called ``constant shifts'' must be separated. In ordinary
quantum mechanics, $H\mapsto H+CI$ with the stress tensor held fixed
changes only an overall phase. A real vacuum-density difference changes
the stress tensor and enters the semiclassical source equation. The first
is nonphysical bookkeeping; the second remains physical even if a chosen
interface cannot see it. The distinction is standard, not a new
cosmological-constant cancellation mechanism.

The appropriate operational rule is therefore a \emph{source-completion
principle}: before predicting a response from restricted data, either
certify the needed source observables or retain their residual uncertainty.
For finite sectors, Theorem~\ref{thm:minimax} gives the sharp residual
floor and Proposition~\ref{prop:finite} connects it to records. This rule
does not prescribe a new gravitational coupling. A stronger law that
filters a physical source would require its own action, conservation and
backreaction analysis; philosophical closure cannot supply those equations.

\subsection{One material, one calibration}
\label{sec:material-vacuum}

The local Gaussian material in HAFF D supplies a positive comparison
\cite{Liu2026HAFFD}. With a fixed texture, masses and boundaries, its
active spatial form contains
\begin{equation}
 q_a(\phi,\Theta)=\int
 \bigl(|\phi'+a(x)\Theta|^2+|\Theta'|^2+m^2|\Theta|^2\bigr)dx,
 \qquad m>0.
\end{equation}
On a finite Dirichlet interval, with $a$ smooth and compactly supported
away from its boundary, the matched operator is
\begin{equation}
 K_a=\begin{pmatrix}-\partial^2&-\partial\circ a\\
                    a\partial&-\partial^2+a^2+m^2\end{pmatrix}.
\end{equation}
The off-diagonal entries are adjoints. The positive squares establish
stability. Reducing to a minimal operator gives a potential-trace
difference $a^2/2$; the standard heat coefficient
\cite{Vassilevich2003} then yields
\begin{equation}
 \Delta E_\Lambda
 =\frac1{8\pi}\int a(x)^2dx\log(\Lambda/\mu)+C_\mu[a]+o(1).
 \label{eq:vacuumlog}
\end{equation}
Matched boundaries cancel their coefficient differences. HAFF D gives
the reduction and the same material action's one-loop counterterm; the
logarithm is not a universal finite gravitational cost.

Changing $a$ changes the coupled fluctuation operator and a calibrated
spectrum or force response. Once material couplings and the finite
normalization are fixed by a common comparison, the response cannot be
removed by choosing a separate energy zero for each configuration.
An unchanged profile with a new Hamiltonian label is still bookkeeping.
An interface that misses the profile-sensitive observable has a
source-certification problem, not a different physical material.

The counterterm and sensor must be included when varying a common
renormalized action. Under a covariant material-functional and state
prescription, its Ward identity includes material equations and any
compensating dilaton force \cite{Liu2026HAFFD}. A static Euclidean or
in--out vacuum comparison does not implement causal control. A controller
that changes $a$ requires work, retained records and an initial-state
retarded response. Neither \eqref{eq:vacuumlog} nor a visibility label
already supplies that controller.

\section{Visibility without an eligibility shortcut}
\label{sec:visibility}

Useful evidence for a vacuum comparison includes causal accessibility,
boundary conditions, state restriction or modular data, actual probe
response, and independently specified asymptotic or sector labels. These
are a menu of tests to be fixed \emph{before} calculating the response,
not a five-item law derived from Q-RAIF A. A change in one may provide
evidence of a physical difference. An unchanged restricted output
establishes indistinguishability in that menu; it cannot classify all
possible underlying differences as energy conventions.

Casimir configurations change boundaries, modes and stress response
\cite{Casimir1948}; phase transitions can change spectra, condensates and
renormalized vacuum parameters. A trace anomaly is a physical
renormalized stress effect rather than a phase relabeling
\cite{Wald1978}. Each needs a matched source calculation. Their possible
gravitational significance does not depend on passing an observer's
preferred visibility test. Conversely, a counterterm or bare zero-point
sum cannot be assessed independently of the common renormalization
condition that fixes the physical stress.

This locates the difference from proposed modifications of gravity.
Unimodular or trace-free formulations change the role of the trace
equation; sequestering introduces constraints on vacuum contributions;
degravitation changes the long-wavelength response
\cite{Ellis2011,KaloperPadilla2014,ArkaniHamed2002}. The present source
certificate does none of these. It bounds what data determine about a
specified source, before that source is inserted into a specified
response law. Induced-gravity approaches use matter response to construct
dynamics and must still match the source and its calibration.

The vacuum-curvature test is consequently two-part. First distinguish a
nonphysical offset from a real changed stress using a common material
model and calibration. Then ask whether the declared interface certifies
that stress well enough to discriminate the candidate response laws.
Here source eligibility precedes a numerical response estimate: it means
establishing a physical, commonly calibrated source difference, not
declaring an unobserved source ineligible to gravitate.
A physical vacuum shift in a null blind direction fails the second part
until the interface is completed. Computing curvature and using its
change as the original visibility datum would make the test circular.

This does not explain why physical Higgs, QCD or phase-transition vacuum
contributions produce no curvature at a naive large scale. That remains
the cosmological-constant problem. The source certificate prevents a
spurious solution by access terminology and identifies what a genuine
source-filtering proposal must measure and predict.

\section{Closure and the location of the reader}
\label{sec:closure}

The certificate itself assumes a source observable, a readout channel and
an instrument. Can a more complete theory generate them? The question
has content even if all three can be represented mathematically. A
stored description of an instrument is not its physical execution, and
specifying a channel is not a derivation of the body, coupling and records
that implement it. An internal model can successfully calibrate some of
these components without thereby explaining the origin of every supplied
structure.

\begin{conjecture}[Readout closure]
A complete physical account of an operational spacetime must specify or
dynamically generate the interfaces that select its records, sectors and
source data, including the physical resources that keep them usable.
An interface may not be fully internalizable within the description it
makes available without an additional level of structure.
\end{conjecture}

The first sentence states the explanatory demand; the second is the
stronger possibility to be tested. A typed fixed point, stratified
description or dynamically formed executor could meet that demand.
Finite programming obstructions and the hidden-source theorem do not
rule those routes out. Nor does an unobserved operator imply that no
observer-independent world exists. The conjecture concerns whether the
completed description accounts for the conditions under which its own
distinctions can be physically made.

Multi-observer descriptions impose record-consistency assumptions
\cite{Rovelli1996,FrauchigerRenner2018}; island calculations assign
radiation and interior sectors \cite{Penington2020,Almheiri2019};
cosmological measures require sampling rules. These are different places
to test whether the account includes its physical interface. AdS/CFT,
for example, may provide an internal description of bulk access through
boundary states and operators. Its success would be evidence for closure,
not a reason to redefine the goal.

A useful test must freeze the meaning of the interface before examining
a claimed complete construction. The test should ask which structures
are primitive, which are generated, how records are physically acquired,
and whether source uncertainty and executor costs are accounted for.
If a closed theory supplies these objects internally, that is a positive
closure result and evidence against the stronger sentence of the
conjecture. It cannot be excluded by renaming every successful
construction an external readout. Conversely, leaving the necessary
instrument as an unpriced oracle does not answer the explanatory demand.

There is a philosophical difference between a world seen incompletely
and a world made only by seeing. The hidden-source example insists on
the former: two worlds can differ in energy while the observer receives
the same record. A new interaction can reveal the difference. The
interface both limits knowledge and gives a physical route to improve
it. ``No view from nowhere'' is most useful as a demand to locate that
route inside the theory, rather than as a prohibition on seeking an
observer-independent account \cite{Nagel1986}.

\section{Conclusion}

Readout blindness has a quantitative source meaning. In a finite sector,
the dual-range distance is exactly the unavoidable error of any estimate
from the ideal output state. An actual instrument adds a separate sample
budget. The positive two-qubit construction shows how the missing
interaction can be recovered through a second priced record of the same
source, with the complete source operator reconstructed by the measured
joint scores.

Null stress misses the metric trace, including real vacuum-density
differences. The matched material comparison shows why changing a
coupling cannot be treated as independently resetting an energy zero.
Both results require physical source completion before a restricted
interface can support a curvature prediction. They do not justify a law
in which whatever a reader misses ceases to gravitate.

The closure conjecture remains broader than these theorems. It asks how
the interfaces that certify and act on sources come into being, and
whether a description can account for its own executor. Concrete source
certificates and completion protocols sharpen that question without
closing it by definition. A theory that genuinely generates its records,
source map and usable apparatus would advance the programme even if it
overturned the conjecture's strongest form.

\begin{thebibliography}{99}
\raggedright
\bibitem{Watrous2018} J.~Watrous, \emph{The Theory of Quantum Information}, Cambridge University Press (2018), Chapter 1, Schatten norm duality.
\bibitem{Hoeffding1963} W.~Hoeffding, Probability inequalities for sums of bounded random variables,
\emph{Journal of the American Statistical Association} \textbf{58}(301), 13--30 (1963), doi:10.1080/01621459.1963.10500830.
\bibitem{BretagnolleHuber1978} J.~Bretagnolle and C.~Huber, Estimation des densit\'es: risque minimax,
\emph{S\'eminaire de Probabilit\'es XII}, Lecture Notes in Mathematics \textbf{649}, 342--363 (1978).
\bibitem{Liu2026HAFFD} S.~Liu, \emph{Selection, Maintenance and Gravitational Response: Modular Constraints and a Physical Return}, 2026; series record DOI: 10.5281/zenodo.18388881.
\bibitem{Liu2026PaperE} S.~Liu, \emph{Measurement as Accessibility: Measurement Permissions, Stable Records and Public Readout}, 2026; series record DOI: 10.5281/zenodo.18400065.
\bibitem{Liu2026PaperF} S.~Liu, \emph{Records, Dissipation and Intervention: Conditions for a Temporal Direction}, 2026; series record DOI: 10.5281/zenodo.18400425.
\bibitem{Weinberg1989} S.~Weinberg, The cosmological constant problem,
\emph{Reviews of Modern Physics} \textbf{61}, 1--23 (1989).
\bibitem{Nagel1986} T.~Nagel, \emph{The View from Nowhere}, Oxford University Press (1986).
\bibitem{Wald1977} R.~M.~Wald, The back reaction effect in particle creation in curved spacetime,
\emph{Communications in Mathematical Physics} \textbf{54}, 1--19 (1977).
\bibitem{Wald1978} R.~M.~Wald, Trace anomaly of a conformally invariant quantum field in curved spacetime,
\emph{Physical Review D} \textbf{17}, 1477--1484 (1978).
\bibitem{Vassilevich2003} D.~V.~Vassilevich, Heat kernel expansion: user's manual,
\emph{Physics Reports} \textbf{388}, 279--360 (2003), arXiv:hep-th/0306138.
\bibitem{Casimir1948} H.~B.~G.~Casimir, On the attraction between two perfectly conducting plates,
\emph{Proceedings of the Koninklijke Nederlandse Akademie van Wetenschappen} \textbf{51}, 793--795 (1948).
\bibitem{Rovelli1996} C.~Rovelli, Relational quantum mechanics,
\emph{International Journal of Theoretical Physics} \textbf{35}, 1637--1678 (1996).
\bibitem{FrauchigerRenner2018} D.~Frauchiger and R.~Renner, Quantum theory cannot consistently describe the use of itself,
\emph{Nature Communications} \textbf{9}, 3711 (2018).
\bibitem{Penington2020} G.~Penington, Entanglement wedge reconstruction and the information paradox,
\emph{Journal of High Energy Physics} \textbf{09}, 002 (2020).
\bibitem{Almheiri2019} A.~Almheiri, N.~Engelhardt, D.~Marolf and H.~Maxfield,
The entropy of bulk quantum fields and the entanglement wedge of an evaporating black hole,
\emph{Journal of High Energy Physics} \textbf{12}, 063 (2019).
\bibitem{Ellis2011} G.~F.~R.~Ellis, H.~van Elst, J.~Murugan and J.-P.~Uzan,
On the trace-free Einstein equations as a viable alternative to general relativity,
\emph{Classical and Quantum Gravity} \textbf{28}, 225007 (2011).
\bibitem{KaloperPadilla2014} N.~Kaloper and A.~Padilla, Sequestering the standard model vacuum energy,
\emph{Physical Review Letters} \textbf{112}, 091304 (2014).
\bibitem{ArkaniHamed2002} N.~Arkani-Hamed, S.~Dimopoulos, G.~Dvali and G.~Gabadadze,
Non-local modification of gravity and the cosmological constant problem, arXiv:hep-th/0209227.
\bibitem{Liu2026QRAIFA} S.~Liu, \emph{Lorentzian Read-out from Local Unitary Propagation: Clifford Relations, Causal Cones and Spectral Mass}, 2026; series record DOI: 10.5281/zenodo.18525876.
\bibitem{Liu2026NullCluster} S.~Liu, \emph{The Null Cluster as Structural Imprint: A Closed Source Task and the Null-Boundary Programme}, 2026; series record DOI: 10.5281/zenodo.20283057.
\end{thebibliography}
\end{document}
